\section{Experimentos e demonstração}

\subsection{Setup experimental}

O experimento utilizou uma base derivada do dataset \nolinkurl{fabriciosantana/discursos-senado-legislatura-56}, contendo 15.729 registros de entrada, dos quais 15.726 correspondiam a discursos efetivamente escritos e 3 foram descartados por ausência de texto útil. A partir desses registros, foram gerados 23.806 chunks, organizados em 120 arquivos Markdown de ingestão. Esses números foram registrados em \nolinkurl{knowledge_openwebui/build_metadata.json}, permitindo rastreabilidade quantitativa da base construída.

Na etapa de preparação da base, adotaram-se como parâmetros principais \texttt{max\_words = 850}, \texttt{overlap\_words = 150} e \texttt{chunks\_per\_file = 200}, buscando equilibrar granularidade e preservação de contexto. Já na configuração operacional validada no Open WebUI, foram definidos \textit{Chunk Size} = 6000, \textit{Chunk Overlap} = 500 e \textit{Chunk Min Size Target} = 2000, com o \textit{Markdown Header Text Splitter} habilitado, uso do modelo de embedding \texttt{text-embedding-3-small} e lote de embeddings de tamanho 32.

O prompt de RAG foi evoluído iterativamente para atender a requisitos centrais do problema institucional: privilegiar o uso prioritário do contexto recuperado, proibir explicitamente a invenção de fatos, exigir declaração de insuficiência quando a informação solicitada não estivesse presente no contexto e incentivar a distribuição de citações ao longo da resposta, complementada, quando pertinente, por uma seção final de referências. Tal configuração alinha-se à literatura que enfatiza explicabilidade, auditabilidade e governança em sistemas RAG \parencite{deene2025raggdpr,said2025ragPractice}.

Para a avaliação, o script \texttt{run\_rag\_eval.py} realizou consultas automatizadas à \textit{knowledge base} do Open WebUI via API, gerando saídas em \texttt{.jsonl}, \texttt{.md} e \texttt{.csv}. Esse arranjo permitiu persistência estruturada, revisão qualitativa e comparação manual entre execuções, aproximando o ambiente experimental de práticas de observabilidade e versionamento recomendadas para sistemas RAG em produção.

\subsection{Resultados alcançados}

A avaliação foi conduzida em duas frentes complementares: (i) verificação operacional da base e do pipeline de ingestão e (ii) execução de consultas automatizadas com persistência de artefatos de análise. Na frente operacional, a ingestão dos 120 arquivos foi concluída sem falhas e manteve coerência com os totais registrados em metadados (15.726 discursos úteis e 23.806 chunks), o que confirma a integridade mínima da base para uso em recuperação semântica.

Na frente de consultas, o script \texttt{run\_rag\_eval.py} gerou saídas em \texttt{.jsonl}, \texttt{.md} e \texttt{.csv}, permitindo auditar cada interação (pergunta, resposta e referências retornadas) e comparar execuções com o mesmo conjunto de dados. Esse desenho experimental tornou possível observar padrões recorrentes em vez de depender apenas de impressão subjetiva.

Os achados qualitativos foram consistentes ao longo das execuções analisadas:
\begin{itemize}
  \item \textbf{Consultas factuais e delimitadas} apresentaram maior precisão e melhor aderência ao contexto recuperado.
  \item \textbf{Consultas temáticas moderadamente abertas} produziram respostas úteis, porém com maior variação de completude conforme a formulação da pergunta.
  \item \textbf{Consultas amplas, comparativas ou multifocais} aumentaram a incidência de sínteses genéricas, com maior necessidade de refinamento iterativo da consulta.
\end{itemize}

Esse comportamento é compatível com a literatura: em RAG, o desempenho final depende da interação entre segmentação, estratégia de recuperação, montagem de contexto e redação do prompt, não apenas do modelo gerador isoladamente \parencite{gao2023ragSurvey,singh2025chunkingTypes,said2025ragPractice}.

Também se observou viabilidade prática de respostas com referências explícitas aos trechos recuperados, aspecto central para auditabilidade no domínio público. Quando o contexto disponível era insuficiente, o prompt instruiu o sistema a declarar limitação informacional em vez de completar lacunas por inferência livre, reduzindo risco de alucinação factual.

Em síntese, os resultados sustentam a viabilidade técnica da prova de conceito e evidenciam limites de escopo já esperados para consultas muito amplas. Como próximo passo metodológico, recomenda-se consolidar uma matriz de avaliação com conjunto fixo de perguntas, rubrica explícita e comparação sistemática entre tipos de consulta, conforme plano descrito em \texttt{NOTAS\_REVISAO.md}.

\subsection{Impacto na administração pública}

O impacto potencial do chatbot legislativo com RAG na administração pública pode ser analisado em pelo menos quatro dimensões. Em primeiro lugar, há impacto sobre a transparência, na medida em que as respostas passam a ser ancoradas em discursos públicos e citáveis, facilitando a verificação por terceiros. Em segundo lugar, há ganho de eficiência, uma vez que o sistema reduz o tempo necessário para localizar informações relevantes em grandes acervos documentais. Em terceiro lugar, a solução fortalece a memória institucional, ao facilitar a recuperação de posicionamentos anteriores, temas recorrentes e evidências documentais. Por fim, observa-se potencial de apoio à decisão, na medida em que gabinetes e equipes técnicas passam a dispor de camada adicional de mediação informacional para explorar o acervo legislativo de forma mais estruturada \parencite{martelloteViola2025organizacaoConhecimento,deene2025raggdpr}.

Por operar sobre fontes controladas e permitir referência explícita aos discursos da 56\textordfeminine{} Legislatura, a solução mostra-se mais compatível com exigências de \textit{accountability} do que o uso isolado de um chatbot generalista. Essa característica assume especial relevância em contextos públicos, nos quais a confiabilidade da informação não pode ser dissociada da governança das fontes e da possibilidade de auditoria posterior.
